\section{Discussion}\label{sec:discussion}

The seven-model demonstration in \Cref{sec:demonstration} has two readings, depending on which audience is doing the reading.

\subsection{For practitioners deploying volatility forecasts}

The practical takeaway is that the choice between top-tier models in a typical volatility-forecasting comparison is often statistically indistinguishable. A risk manager comparing a GJR-GARCH against a LightGBM against an equal-weight ensemble on S\&P 500 data is not, in expectation, choosing between forecasters of different accuracy; she is choosing between forecasters of statistically equal accuracy with different operational properties (parameter count, retraining cost, interpretability, regulator-defensibility under SR 11-7 / OCC 2011-12 / EU AI Act). The honest framing of the choice is therefore not ``which is most accurate?'' but ``which is most defensible?'' \citep{FedSR117, OCC2011, EUAIAct}.

A second practical takeaway, less obvious from the original paper but visible in our \Cref{tab:demo_dm}: the differences that \emph{are} statistically significant tend to be those that involve a clearly inferior model (HAR-RV losing to anything with a leverage-effect feature; GARCH losing to GJR-GARCH because the asymmetry term matters in this test period). The differences between top-tier models are not significant. This is consistent with the broader empirical finding from \citet{KhanVol2026} that the ranking of top models is regime-dependent and unstable across subperiods. The implication is that an ensemble has a robustness argument even when it does not have an accuracy argument: it averages the regime-dependent failures of its components.

\subsection{For the methodological literature}

The takeaway here is that single-number ranking tables in published volatility-forecasting papers are systematically over-interpreted. Of 21 pairwise comparisons in our demonstration, only 6 (29\%) are significant at 5\%. The standard practice of reporting the lowest-loss model as the ``winner'' implies a discriminating power that the underlying statistics do not support. We would expect the same to be true of most published comparisons in this subfield; the seven-model panel of \citet{KhanVol2026} is unusual mainly in being unusually well-documented.

Three editorial recommendations follow.

\textit{Default to reporting MCS survivors instead of single rankings.} A 90\% MCS that keeps six of seven models is informative content. A ranking table with bold values implying a clear winner where no clear winner exists is misinformation, even when each individual cell is correct. The MCS framing makes the underlying ambiguity visible.

\textit{Report Diebold-Mariano significance for any pairwise comparison the prose makes.} If the prose says ``model X beats model Y on QLIKE,'' the prose should also say at what p-value. If the answer is ``not significantly,'' the prose should say so plainly. \texttt{vol-eval}'s \texttt{dm\_test} returns the p-value as a one-line addition to any existing comparison.

\textit{Be especially careful with multi-model comparisons.} The DM test is for two models; using it as a battery of pairwise comparisons across many models without multiplicity adjustment will inflate the family-wise error rate and produce illusory ``wins.'' The MCS handles this correctly. The SPA / Reality Check handle the related ``one benchmark vs.\ many competitors'' question.

\subsection{The queued empirical extension}

The seven-model demonstration is one paper. The interesting question is how the findings generalize. Specifically: of the published volatility-forecasting comparison papers from the last five years, what fraction would survive a re-analysis under \texttt{vol-eval}?

We have scaffolded the empirical extension in the project repository. Phase 2 of the project, scheduled for completion by mid-July 2026, will:

\begin{enumerate}[leftmargin=*]
    \item Identify 30 published volatility-forecasting comparison papers from 2018--2025 via a Crossref + OpenAlex + arXiv bibliographic search, filtered for horse-race language (papers comparing $\geq 2$ model families and reporting out-of-sample loss evaluation). After narrowing and prioritising papers with openly-accessible PDFs, the final sample at the time of writing comprises six papers from the \emph{Journal of Financial Econometrics}, thirteen from three open-access MDPI venues (\emph{Mathematics}, \emph{Risks}, \emph{Journal of Risk and Financial Management}), and eleven arXiv q-fin preprints. The originally-targeted Wiley and Elsevier journals (\emph{Journal of Forecasting}, \emph{International Journal of Forecasting}, \emph{Journal of Empirical Finance}, \emph{Quantitative Finance}, \emph{Journal of Financial Data Science}) are absent from this analysis-stage sample due to PDF access constraints during programmatic download; we record this access-driven sample-frame deviation as a Phase 2 limitation, noting that the open-access MDPI venues we use are indexed in DOAJ and apply transparent (if briefer than tier-1) peer review.
    \item Apply the Reproducibility Disclosure Score rubric of \citet{Khan2026Survey} to each, scoring 0--2 based on code and data availability.
    \item For papers that score 2 (full reproducibility), clone the code, rerun, extract the per-day forecast time series for each model, and apply the full \texttt{vol-eval} battery (DM, MCS, SPA) to the author-claimed pairwise winner against the runner-up.
    \item For papers that score 1 (public data, no code), re-implement the headline models from the paper's description and apply the same battery.
    \item For papers that score 0 (vendor data, no code), document the unreproducibility in the disclosure tabulation and skip the significance analysis.
    \item Tabulate, across the reproducible subset, what fraction of headline ``wins'' survive (i) DM at 5\%, (ii) MCS at 90\%, (iii) SPA-consistent at 5\%.
\end{enumerate}

We do not pre-announce a specific predicted finding for the empirical extension because the sample is yet to be collected. Based on the seven-model demonstration here and on the editorial-survey evidence in \citet{Khan2026Survey}, our prior is that a substantial minority (perhaps 30--50\% by DM, more by MCS) of headline ``wins'' will not survive significance testing. The point of the empirical extension is to test this prior with data, not to assume it.

\subsection{Limitations}

Three limitations of the present paper deserve noting.

First, the demonstration is on a single seven-model panel. The findings about MCS survival, SPA non-rejection, and the failure of the ensemble vs.\ GJR-GARCH ``win'' to be significant are about that specific panel and that specific test period. The same demonstration on a different asset class, a different forecast horizon, or a different test window would produce different specific numbers. The methodological conclusion (that significance testing matters) generalizes; the specific p-values do not.

Second, \texttt{vol-eval}'s implementation choices for the bootstrap and for the SPA centering follow Hansen (2005) and the \texttt{arch} package's conventions. Alternative bootstrap designs (block bootstrap with fixed block length, parametric bootstrap, GARCH-type bootstrap) would produce slightly different p-values. We default to the stationary bootstrap because it is the standard in the published literature, but users who want a different design can pass it as an argument. The package's API does not foreclose alternatives; the current release simply ships one default.

Third, the paper is a methods paper with a self-contained empirical demonstration, not a re-analysis of the published literature. The queued Phase 2 extension is the project that produces the broader empirical claim. We have been careful not to overclaim from the seven-model demonstration; the tabular findings should be cited as such, and the broader empirical question should be cited as ``in production'' until Phase 2 ships.
